\section{Related Work}
\paragraph{Label and category representation.}
NLP models often use labels as more than output IDs. Labels can be embedded from their names, expanded with descriptions, matched to examples, or treated as prompts. Weakly supervised and zero-shot classifiers use label semantics to connect text instances to categories, while prototype-based methods represent a class by one or more vectors estimated from examples \citep{snell2017prototypical,reimers2019sentencebert,meng2020lotclass}. Our setting is closest to example-derived category representation: the label vector is induced from constituent members rather than from the label name alone. The distinction matters because a member-derived representation can be evaluated both as a category prototype and as a downstream model component.

\paragraph{Similarity and downstream validity.}
Embedding similarity is not a single universal relation. Two texts may be similar by topic, reasoning step, difficulty, prerequisite relation, or expected student response pattern. Recent evaluation work in representation learning emphasizes that similarity should be tied to the target relation being modeled rather than treated as a generic property. Our study operationalizes that point: KC prototypes are tested against category membership and label-name structure, then separately against KT-relevant functional graphs and downstream DKT performance.

\paragraph{Educational text and knowledge tracing.}
Educational NLP has shown that question text can support tasks such as difficulty prediction, content tagging, and feedback generation. KT models, however, are trained on student interaction sequences and depend on temporal and behavioral regularities as much as text content. DKT \citep{piech2015dkt} represents each KC-response pair as an input symbol whose embedding is learned for prediction. Chapter 5 of the thesis replaces this learned initialization with text-derived KC-response embeddings and finds little robust improvement. Chapter 6 studies a different problem, redefining or reclustering the KC space itself; that direction is treated here only as separate work, not as evidence for the present paper.
